// ============================================================
//  calculatrice.tex — Mini calculatrice graphique pour TetraOS
//
//  Utilise : include @package.AppCore
//
//  Fonctionnalités :
//    - Chiffres 0–9, virgule décimale
//    - Opérations : + - * /
//    - [=]   calculer le résultat
//    - [C]   tout effacer
//    - [Del] effacer le dernier caractère
//    - Affichage : expression en bleu, résultat en vert (ou rouge si erreur)
//    - Après [=], la prochaine saisie repart de zéro
//
//  Note : le moteur TEX évalue les expressions arithmétiques
//  directement (3+4*2, 10/4, etc.) via tex_eval_expr, donc
//  on passe la chaîne expr directement au moteur pour le calcul.
// ============================================================

include @package.AppCore

// ── Variables d'état ─────────────────────────────────────────
var expr     = "0"     // Expression en cours
var op_faite = 0       // 1 si on vient de faire = (prochaine saisie repart de zéro)
var MAX_LEN  = 24      // Longueur max de l'expression

// ── Fenêtre ───────────────────────────────────────────────────
app.init()
app.window("Calculatrice", 300, 180, 314, 410, win)

// ── Zone d'affichage ─────────────────────────────────────────

// Label expression (ce que l'utilisateur tape)
app.label(win, 12, 12, "0", lbl_expr)
app.label_color(lbl_expr, 0x0077AACC)

// Label résultat (affiché après =)
app.label(win, 12, 40, "", lbl_result)
app.label_color(lbl_result, 0x00AAFFAA)

// Ligne de séparation décorative
app.label(win, 12, 68, "______________________________", lbl_sep)
app.label_color(lbl_sep, 0x00223344)

// ── Grille de boutons (4×5) ───────────────────────────────────
//
//  Col :   10     80    150    220
//  Lg1:  [ C  ] [Del] [ / ] [ * ]    y=82
//  Lg2:  [ 7  ] [ 8 ] [ 9 ] [ - ]    y=118
//  Lg3:  [ 4  ] [ 5 ] [ 6 ] [ + ]    y=154
//  Lg4:  [ 1  ] [ 2 ] [ 3 ] [   ]    y=190   \
//  Lg5:  [  0   ] [ . ]     [ = ]    y=226   /  = sur 2 lignes

// Ligne 1
app.btn(win,  10, 82, 64, 32, "C",   b_c)
app.btn(win,  80, 82, 64, 32, "Del", b_del)
app.btn(win, 150, 82, 64, 32, "/",   b_div)
app.btn(win, 220, 82, 64, 32, "*",   b_mul)

// Ligne 2
app.btn(win,  10, 118, 64, 32, "7", b_7)
app.btn(win,  80, 118, 64, 32, "8", b_8)
app.btn(win, 150, 118, 64, 32, "9", b_9)
app.btn(win, 220, 118, 64, 32, "-", b_sub)

// Ligne 3
app.btn(win,  10, 154, 64, 32, "4", b_4)
app.btn(win,  80, 154, 64, 32, "5", b_5)
app.btn(win, 150, 154, 64, 32, "6", b_6)
app.btn(win, 220, 154, 64, 32, "+", b_add)

// Ligne 4
app.btn(win,  10, 190, 64, 32, "1", b_1)
app.btn(win,  80, 190, 64, 32, "2", b_2)
app.btn(win, 150, 190, 64, 32, "3", b_3)
app.btn(win, 220, 190, 64, 68, "=", b_eq)    // = sur 2 lignes (hauteur 68)

// Ligne 5
app.btn(win,  10, 226, 134, 32, "0", b_0)    // 0 large (2 colonnes)
app.btn(win, 150, 226,  64, 32, ".", b_dot)

// ── Fonctions TEX ─────────────────────────────────────────────

// Ajoute un caractère à l'expression
func ajouter(c)
    // Si on vient de faire =, recommencer à zéro
    if op_faite == 1 {
        var expr     = c
        var op_faite = 0
        app.set_label(lbl_result, "")
        app.set_label(lbl_expr, expr)
        app.label_color(lbl_expr, 0x0077AACC)
        return
    }
    // Remplacer le "0" initial par le premier chiffre
    if expr == "0" {
        if c == "+" || c == "-" || c == "*" || c == "/" || c == "." {
            var expr = "0" + c
        } else {
            var expr = c
        }
    } else {
        // Limiter la longueur de l'expression
        var n = str.len(expr)
        if n < MAX_LEN {
            var expr = expr + c
        }
    }
    app.set_label(lbl_expr, expr)
end

// Efface le dernier caractère
func effacer_dernier()
    var n = str.len(expr)
    if n > 1 {
        var expr = str.sub(expr, 0, n - 1)
    } else {
        var expr = "0"
    }
    app.set_label(lbl_expr, expr)
    app.label_color(lbl_expr, 0x0077AACC)
end

// Tout effacer
func tout_effacer()
    var expr     = "0"
    var op_faite = 0
    app.set_label(lbl_expr,   "0")
    app.set_label(lbl_result, "")
    app.label_color(lbl_expr,   0x0077AACC)
    app.label_color(lbl_result, 0x00AAFFAA)
end

// Calcule et affiche le résultat
// Le moteur TEX évalue les expressions arithmétiques directement.
// On passe expr au moteur via une assignation dynamique :
//   var res = <valeur de expr évaluée comme expression>
// Comme TEX évalue le membre droit comme expression arithmétique,
// on doit passer expr comme une expression à évaluer.
// Astuce : on évalue expr en tant qu'expression TEX via math.eval.
func calculer()
    // Evaluer l'expression arithmétique contenue dans expr
    // math.eval(var) lit la valeur string de var et l'évalue comme expression
    var res = math.eval(expr)

    // Détecter la division par zéro dans l'expression
    // (le moteur retourne 0 ou inf, on vérifie via str.find)
    var pos_div0 = str.find(expr, "/0")

    if pos_div0 >= 0 {
        // Vérifier que c'est bien "/0" seul et pas "/0.5" etc.
        var apres = str.sub(expr, pos_div0 + 2, pos_div0 + 3)
        if apres == "" || apres == "+" || apres == "-" || apres == "*" || apres == "/" {
            app.set_label(lbl_result, "Erreur : /0")
            app.label_color(lbl_result, 0x00FF4444)
            app.label_color(lbl_expr,   0x00884444)
            var op_faite = 1
            return
        }
    }

    // Convertir le résultat en chaîne et afficher
    var res_str = str(res)
    app.set_label(lbl_result, "= " + res_str)
    app.label_color(lbl_result, 0x00AAFFAA)
    app.label_color(lbl_expr,   0x00445566)   // griser l'expression

    var op_faite = 1
end

// ── Boucle principale ─────────────────────────────────────────
var run = 1
while run == 1 {
    app.tick()

    // ── Chiffres ──────────────────────────────────────────────
    app.btn_clicked(b_0, v) if v == 1 { ajouter("0") }
    app.btn_clicked(b_1, v) if v == 1 { ajouter("1") }
    app.btn_clicked(b_2, v) if v == 1 { ajouter("2") }
    app.btn_clicked(b_3, v) if v == 1 { ajouter("3") }
    app.btn_clicked(b_4, v) if v == 1 { ajouter("4") }
    app.btn_clicked(b_5, v) if v == 1 { ajouter("5") }
    app.btn_clicked(b_6, v) if v == 1 { ajouter("6") }
    app.btn_clicked(b_7, v) if v == 1 { ajouter("7") }
    app.btn_clicked(b_8, v) if v == 1 { ajouter("8") }
    app.btn_clicked(b_9, v) if v == 1 { ajouter("9") }

    // ── Point décimal ─────────────────────────────────────────
    app.btn_clicked(b_dot, v) if v == 1 { ajouter(".") }

    // ── Opérateurs ────────────────────────────────────────────
    app.btn_clicked(b_add, v) if v == 1 { ajouter("+") }
    app.btn_clicked(b_sub, v) if v == 1 { ajouter("-") }
    app.btn_clicked(b_mul, v) if v == 1 { ajouter("*") }
    app.btn_clicked(b_div, v) if v == 1 { ajouter("/") }

    // ── Actions ───────────────────────────────────────────────
    app.btn_clicked(b_del, v) if v == 1 { effacer_dernier() }
    app.btn_clicked(b_c,   v) if v == 1 { tout_effacer()    }
    app.btn_clicked(b_eq,  v) if v == 1 { calculer()        }

    // Vérifier si la fenêtre est encore ouverte
    app.running(run)
}
